% !TeX root = ../../main.tex

\section{Introducción}

MapReduce es un modelo de programación utilizado para procesar grandes
cantidades de datos de manera distribuida. Su funcionamiento consiste en
dividir un problema en tareas pequeñas que pueden ejecutarse al mismo tiempo
en diferentes computadoras.

Este modelo forma parte de Apache Hadoop y se utiliza principalmente para
trabajos por lotes.

\section{¿Qué es MapReduce?}

MapReduce es un modelo de programación y procesamiento distribuido que permite
analizar grandes volúmenes de información utilizando varias computadoras de un
clúster.

Los datos se representan mediante pares de clave y valor. El procesamiento se
divide principalmente en dos funciones:

\begin{itemize}
    \item \textbf{Map:} procesa partes del conjunto de datos y genera resultados
    intermedios.
    \item \textbf{Reduce:} combina los resultados intermedios para obtener un
    resultado final.
\end{itemize}

El framework se encarga de distribuir las tareas, coordinar las computadoras,
transferir los datos y repetir una tarea cuando ocurre una falla.

\section{Antecedentes de MapReduce}

Las ideas de \textit{map} y \textit{reduce} tienen su origen en los lenguajes
de programación funcional. Estas operaciones permiten aplicar una función a
varios elementos y posteriormente reunir sus resultados.

En 2004, Jeffrey Dean y Sanjay Ghemawat presentaron el modelo MapReduce
desarrollado por Google. Su propósito era simplificar el procesamiento de
grandes cantidades de información mediante clústeres formados por miles de
computadoras.

Posteriormente, el proyecto Apache Hadoop creó una implementación de código
abierto basada en este modelo. Esto permitió que otras organizaciones pudieran
utilizar MapReduce junto con HDFS para almacenar y procesar grandes volúmenes
de datos.

\section{Arquitectura de MapReduce}

La arquitectura de MapReduce combina el almacenamiento distribuido de HDFS con
varios componentes encargados de administrar y ejecutar las tareas.

Sus elementos principales son los siguientes:

\begin{itemize}
    \item \textbf{Cliente:} envía el programa, los datos de entrada y la
    configuración del trabajo.

    \item \textbf{HDFS:} almacena los archivos de entrada y los resultados
    finales de manera distribuida.

    \item \textbf{ResourceManager:} administra los recursos disponibles en el
    clúster.

    \item \textbf{ApplicationMaster:} coordina la ejecución de un trabajo
    específico de MapReduce.

    \item \textbf{NodeManager:} administra los recursos de cada nodo y ejecuta
    las tareas asignadas.

    \item \textbf{Mappers:} procesan las divisiones de los datos y producen
    pares de clave y valor intermedios.

    \item \textbf{Reducers:} reciben los valores agrupados por clave y generan
    los resultados finales.
\end{itemize}

En Hadoop 1 se utilizaban los componentes \textit{JobTracker} y
\textit{TaskTracker}. A partir de Hadoop 2, estas funciones fueron separadas y
administradas mediante YARN, utilizando componentes como ResourceManager,
ApplicationMaster y NodeManager.

\section{Principales características de MapReduce}

Las características más importantes de MapReduce son:

\begin{itemize}
    \item Procesa grandes cantidades de datos por lotes.
    \item Divide un trabajo en varias tareas pequeñas.
    \item Ejecuta las tareas de forma paralela y distribuida.
    \item Trabaja con pares de clave y valor.
    \item Puede aumentar su capacidad agregando más nodos al clúster.
    \item Aprovecha la ubicación de los datos para reducir transferencias.
    \item Detecta fallas y puede volver a ejecutar las tareas.
    \item Permite procesar datos estructurados, semiestructurados y no
    estructurados.
    \item Automatiza la distribución, coordinación y supervisión del trabajo.
\end{itemize}

\section{Funciones Map() y Reduce()}

\subsection*{Función Map()}

La función \texttt{Map()} recibe un par de clave y valor de entrada. Después
aplica una operación y genera uno o varios pares intermedios.

Su representación general es:

\begin{center}
\texttt{Map(K1, V1)} $\longrightarrow$ \texttt{lista(K2, V2)}
\end{center}

Por ejemplo, en un contador de palabras, la función Map separa una línea de
texto y genera el valor 1 para cada palabra encontrada.

\begin{center}
\texttt{Hadoop procesa datos}
\end{center}

La salida sería:

\begin{center}
\texttt{(Hadoop, 1), (procesa, 1), (datos, 1)}
\end{center}

\subsection*{Función Reduce()}

La función \texttt{Reduce()} recibe una clave y la lista de valores
relacionados con ella. Su trabajo consiste en combinar o resumir esos valores.

Su representación general es:

\begin{center}
\texttt{Reduce(K2, lista(V2))} $\longrightarrow$ \texttt{(K3, V3)}
\end{center}

Por ejemplo, si recibe:

\begin{center}
\texttt{(Hadoop, [1, 1, 1])}
\end{center}

La función suma los valores y produce:

\begin{center}
\texttt{(Hadoop, 3)}
\end{center}

\section{Flujo de procesamiento}

El procesamiento de datos con MapReduce se realiza mediante las siguientes
etapas:

\begin{enumerate}
    \item \textbf{Entrada:} los archivos de entrada se almacenan normalmente
    en HDFS.

    \item \textbf{División:} los archivos se dividen en fragmentos llamados
    \textit{input splits}. Cada fragmento será procesado por una tarea Map.

    \item \textbf{Lectura:} el RecordReader convierte los registros en pares de
    clave y valor que puede recibir el Mapper.

    \item \textbf{Map:} cada Mapper procesa su fragmento y genera pares
    intermedios.

    \item \textbf{Combiner:} de manera opcional, combina resultados locales
    para reducir la cantidad de información enviada por la red.

    \item \textbf{Partition:} determina a qué Reducer se enviará cada clave.

    \item \textbf{Shuffle y Sort:} los resultados se transfieren, ordenan y
    agrupan de acuerdo con sus claves.

    \item \textbf{Reduce:} cada Reducer procesa una clave y su lista de valores.

    \item \textbf{Salida:} los resultados finales se guardan nuevamente en
    HDFS.
\end{enumerate}

De forma resumida, el flujo es:

\begin{center}
Entrada $\longrightarrow$ División $\longrightarrow$ Map
$\longrightarrow$ Shuffle y Sort $\longrightarrow$ Reduce
$\longrightarrow$ Salida
\end{center}

\section{Ventajas y desventajas}

\subsection*{Ventajas}

\begin{itemize}
    \item Permite procesar terabytes o petabytes de información.
    \item Ejecuta varias tareas al mismo tiempo.
    \item Puede funcionar en computadoras de costo relativamente bajo.
    \item Es escalable porque se pueden agregar más nodos.
    \item Ofrece tolerancia a fallas mediante la repetición de tareas.
    \item Reduce movimientos innecesarios al procesar los datos cerca de donde
    están almacenados.
    \item Simplifica algunos trabajos de procesamiento distribuido.
\end{itemize}

\subsection*{Desventajas}

\begin{itemize}
    \item Está orientado al procesamiento por lotes y no al tiempo real.
    \item Puede presentar una latencia alta.
    \item Las operaciones de lectura y escritura en disco pueden hacerlo lento.
    \item La transferencia de datos durante \textit{shuffle} consume recursos
    de la red.
    \item No es eficiente para conjuntos de datos pequeños.
    \item Los algoritmos iterativos pueden requerir varios trabajos MapReduce.
    \item La fase Reduce normalmente debe esperar a que las tareas Map terminen.
\end{itemize}

\section{Ejemplo de uso: contador de palabras}

Un ejemplo sencillo de MapReduce consiste en contar cuántas veces aparece cada
palabra en un archivo.

Supongamos que el archivo contiene las siguientes líneas:

\begin{center}
\texttt{Hadoop procesa datos}\\
\texttt{Hadoop almacena datos}
\end{center}

\subsection*{Resultado de Map}

Cada Mapper separa las palabras y genera los siguientes pares:

\begin{itemize}
    \item \texttt{(Hadoop, 1)}
    \item \texttt{(procesa, 1)}
    \item \texttt{(datos, 1)}
    \item \texttt{(Hadoop, 1)}
    \item \texttt{(almacena, 1)}
    \item \texttt{(datos, 1)}
\end{itemize}

\subsection*{Resultado de Shuffle y Sort}

El framework agrupa los valores que tienen la misma clave:

\begin{itemize}
    \item \texttt{(Hadoop, [1, 1])}
    \item \texttt{(datos, [1, 1])}
    \item \texttt{(procesa, [1])}
    \item \texttt{(almacena, [1])}
\end{itemize}

\subsection*{Resultado de Reduce}

Finalmente, la función Reduce suma los valores de cada palabra:

\begin{itemize}
    \item \texttt{(Hadoop, 2)}
    \item \texttt{(datos, 2)}
    \item \texttt{(procesa, 1)}
    \item \texttt{(almacena, 1)}
\end{itemize}

Este mismo procedimiento puede aplicarse a archivos mucho más grandes,
distribuyendo el trabajo entre diferentes nodos del clúster.

\section{Conclusión}

MapReduce permite procesar grandes volúmenes de información mediante la
división del trabajo en tareas Map y Reduce. Sus principales ventajas son el
procesamiento paralelo, la escalabilidad y la tolerancia a fallas.

Aunque no es la mejor opción para aplicaciones en tiempo real o procesos
iterativos, continúa siendo un modelo importante para comprender el
funcionamiento del procesamiento distribuido dentro del ecosistema Hadoop.